\documentclass[11pt,a4paper]{article}
\usepackage[T1]{fontenc}
\usepackage[utf8]{inputenc}
\usepackage[main=american,finnish]{babel}
\usepackage{lmodern}
\usepackage{microtype}
\usepackage[margin=22mm,headheight=44pt,headsep=8mm,footskip=12mm]{geometry}
\usepackage{graphicx}
\usepackage[dvipsnames]{xcolor}
\usepackage{booktabs,tabularx,array,longtable}
\usepackage{enumitem}
\usepackage{amsmath,amssymb}
\usepackage{caption}
\usepackage{float}
\usepackage{fancyhdr}
\usepackage[most]{tcolorbox}
\usepackage[hidelinks]{hyperref}
\usepackage{url}
% Visible amendment mode for the ethics-committee supplement.
% Set \IRBchangesfalse to prepare a clean submission copy.
\usepackage[normalem]{ulem}
\usepackage{soul}
\sethlcolor{yellow}
\newif\ifIRBchanges
\IRBchangestrue
\newcommand{\irbadd}[1]{\ifIRBchanges\hl{#1}\else#1\fi}
\newcommand{\irbdel}[1]{\ifIRBchanges{\color{red!70!black}\sout{#1}}\fi}
\newcommand{\irbreplace}[2]{\ifIRBchanges\irbdel{#1}\allowbreak\,\irbadd{#2}\else#2\fi}


\definecolor{UTUBlue}{HTML}{005F86}
\definecolor{UTULight}{HTML}{EAF3F7}
\definecolor{UTUDark}{HTML}{263746}
\definecolor{SoftGray}{HTML}{F3F5F6}
\definecolor{WarnGold}{HTML}{F6E8B1}
\newcommand{\StudyTitle}{To Renounce Liberty? A Public-Goods Experiment on (Un)Democratic Reversal}
\newcommand{\StudyTitleFI}{Luopuako vapaudesta? Julkishyödykekoe (epä)demokraattisesta suunnanmuutoksesta}
\newcommand{\StudyPI}{Janne Tukiainen}
\newcommand{\StudyEmail}{\href{mailto:janne.tukiainen@utu.fi}{janne.tukiainen@utu.fi}}
\newcommand{\StudyPhone}{+358 50 308 3620}
\newcommand{\DocumentDate}{20 October 2026}
\newcommand{\DocumentVersion}{Version 2.6 -- proposed amendment}
\newcommand{\checkbox}{$\square$}

\pagestyle{fancy}
\fancyhf{}
\fancyhead[L]{\includegraphics[height=13mm]{utu_logo.pdf}}
\fancyhead[R]{\footnotesize\color{UTUDark}\StudyTitle}
\fancyfoot[L]{\footnotesize\color{UTUDark}\DocumentVersion\ -- \DocumentDate}
\fancyfoot[R]{\footnotesize\color{UTUDark}\thepage}
\renewcommand{\headrulewidth}{0.35pt}
\renewcommand{\footrulewidth}{0.25pt}
\setlength{\parindent}{0pt}
\setlength{\parskip}{0.55em}
\setlength{\emergencystretch}{3em}
\setlist{nosep,leftmargin=1.35em,topsep=0.25em}
\renewcommand{\arraystretch}{1.18}
\urlstyle{same}
\hypersetup{colorlinks=true,linkcolor=UTUBlue,urlcolor=UTUBlue,citecolor=UTUBlue,pdfauthor={Janne Tukiainen},pdftitle={\StudyTitle}}
\captionsetup{font=small,labelfont=bf}

\newcommand{\IRBTitle}[2]{%
  \thispagestyle{fancy}
  {\small\color{UTUBlue}\textbf{#1}}\par
  \vspace{0.35em}
  {\LARGE\bfseries\color{UTUDark}#2}\par
  \vspace{0.35em}
  {\large\StudyTitle}\par
  \vspace{0.45em}
  {\small\StudyPI\ $\vert$ \StudyEmail\ $\vert$ \DocumentVersion\ $\vert$ \DocumentDate}\par
  \vspace{0.65em}\hrule\vspace{0.8em}
}

\newcommand{\IRBTitleFI}[2]{%
  \thispagestyle{fancy}
  {\small\color{UTUBlue}\textbf{#1}}\par
  \vspace{0.35em}
  {\LARGE\bfseries\color{UTUDark}#2}\par
  \vspace{0.35em}
  {\large\StudyTitleFI}\par
  \vspace{0.45em}
  {\small\StudyPI\ $\vert$ \StudyEmail\ $\vert$ \DocumentVersion\ $\vert$ \DocumentDate}\par
  \vspace{0.65em}\hrule\vspace{0.8em}
}

\newtcolorbox{infobox}[1][Key information]{
  colback=UTULight,colframe=UTUBlue,coltitle=white,
  fonttitle=\bfseries,title=#1,boxrule=0.7pt,arc=1.5mm,
  left=7pt,right=7pt,top=6pt,bottom=6pt
}
\newtcolorbox{reviewbox}[1][Confirmation required before submission]{
  colback=WarnGold!55,colframe=Orange!75!black,coltitle=black,
  fonttitle=\bfseries,title=#1,boxrule=0.7pt,arc=1.5mm,
  left=7pt,right=7pt,top=6pt,bottom=6pt
}
\newcommand{\sectionbar}[1]{\vspace{0.55em}{\large\bfseries\color{UTUBlue}#1}\par\vspace{-0.15em}}
\newcolumntype{Y}{>{\raggedright\arraybackslash}X}
\newcolumntype{L}[1]{>{\raggedright\arraybackslash}p{#1}}

\begin{document}
\IRBTitle{Attachment 7}{Other Materials Provided to Participants}

This attachment reproduces the substantive participant-facing instructions, checks, expectation items, final questionnaire, and debriefing implemented in oTree. Dynamic values such as the current round, proposal, rate, and payoff are inserted by the software. Research-only solo-testing controls are not shown in production.

\sectionbar{Opening instructions}
The first screen states that registration through \irbreplace{the host laboratory's participant system}{UTU Choice Lab's UTU RPP/TOKA participant system} provides its general rules and privacy information, then gives concise information and requests agreement to this particular study. It identifies Janne Tukiainen, Professor at the Turku School of Economics, University of Turku, as Principal Investigator and contact (\StudyEmail; \StudyPhone), and explains the approximately 60-minute duration, interactive task, payment, minimal foreseeable risks, voluntary participation, coded data, \irbreplace{and removal of any host-laboratory identity-to-code link after payment}{that the detailed research question will be explained after paid decisions, and the separate code-and-amount payment procedure}.

Imagine that you live in a society where public services have been under serious strain for a long time. Parliament has failed to pass a budget for eight months, healthcare waiting times have doubled, and planning for public services has stalled.

During the experiment, you and the other participants are the citizens of this society. You make decisions about a public-services fund. Points you keep count toward your payoff. Points that remain in the fund create group points shared equally among five citizens.

There are \irbreplace{twenty}{fifteen} paid rounds in two blocks\irbreplace{ of ten}{: ten in Block 1 and five in Block 2}. In each round, participants choose whether a selected person's proposal requires group approval or takes effect directly. The method selected by at least three participants is used. The software then selects one of the five citizens uniformly at random, giving each a one-in-five chance in every round.

Under \textbf{Group approval required}, the other four citizens see the complete proposal and at least three must approve. If rejected, everyone chooses their own fund allocation. Under \textbf{Decision takes effect directly}, the proposal is implemented without that approval vote. Results and individual payoffs are then shown. Participants remain in a fixed matching pool of \irbreplace{10 or 15}{5, 10, or 15}\irbreplace{ and are reshuffled into five-person groups after every round}{. Ten- and fifteen-person pools are reshuffled into five-person groups before every round; in a five-person pool, the same people play together throughout}; \irbreplace{identities are never shown}{the software displays no names or identities}.

After all rounds, one round from each block is randomly selected. Points earned in those two rounds determine a performance bonus of up to EUR~5, and a EUR~10 participation payment is added. \irbadd{Your points from these two rounds are converted at a rate of 24 points = EUR~1.} More selected-round points produce a larger bonus. Practice decisions are unpaid. \irbadd{On paid decision screens, a 90-second countdown reminds participants to make their choice and click the confirmation button. At zero, the reminder remains visible and the decision stays open. Questionnaires are untimed.} There is no automatic time limit; the \irbreplace{laboratory manager}{researcher or trained research assistant} monitors progress and may prompt delayed participants.

\sectionbar{Unpaid practice}
Participants practise three decisions:
\begin{enumerate}
  \item \textbf{Make a proposal:} choose how many of each citizen's 20 points should enter the fund and whether up to 20 fund points move to the selected person's payoff. A transfer cannot exceed the amount placed in the fund.
  \item \textbf{Approve or reject:} respond to an example in which every citizen would place 10 points in the fund and 5 fund points would move to the selected person's payoff.
  \item \textbf{Allocation after rejection:} choose an individual allocation after being told that the other four citizens each place 10 points in the fund at the 1.50 rate.
\end{enumerate}

\sectionbar{Comprehension check}
Participants answer all five items correctly before paid play. Incorrect answers return a brief explanation and another attempt.

\begin{enumerate}
  \item \textbf{Who proposes the equal allocation under the two methods?}
  \begin{itemize}
    \item A randomly selected person proposes the allocation under both methods. \emph{[Correct]}
    \item A selected person is used only when the decision takes effect directly.
  \end{itemize}
  \item \textbf{When group approval is required, how many of the four other citizens must approve?}
  \begin{itemize}
    \item At least three of the four other citizens. \emph{[Correct]}
    \item At least two of the four other citizens.
    \item All four other citizens.
  \end{itemize}
  \item \textbf{What happens when a proposal requiring group approval is rejected?}
  \begin{itemize}
    \item Everyone chooses their own fund allocation. \emph{[Correct]}
    \item The proposal takes effect anyway.
  \end{itemize}
  \item \textbf{When a decision takes effect directly, can the other citizens reject the proposal before it is implemented?} Yes / No \emph{[Correct: No]}
  \item \textbf{Which paid round determines your game earnings from Block 1?}
  \begin{itemize}
    \item Every round.
    \item Only the randomly selected round. \emph{[Correct]}
    \item Only Round 10.
  \end{itemize}
\end{enumerate}

\sectionbar{Block 1 information}
Every fund point left after any transfer creates 1.50 group points under Group approval required and 2.50 under Decision takes effect directly. If an approval-required proposal is rejected, participants make individual allocations at the 1.50 rate. The selected person may propose the same allocation and transfer under either method.

\sectionbar{Block 2 updates}
Each matching pool sees one of the following screens after completing Block 1:

\textbf{Conditions improve.} Parliament has passed a budget, healthcare waiting times are falling, and planning for public services has resumed. As part of this recovery, every fund point under Group approval required now creates 2.50 group points. Under Decision takes effect directly, every point left in the fund creates 2.50 group points.

\textbf{Conditions remain under strain.} Parliament has still not passed a budget, healthcare waiting times remain high, and planning for public services remains stalled. Under Group approval required, every fund point creates 1.50 group points. Under Decision takes effect directly, every point left in the fund creates 2.50 group points.

In both updates, the decision rules remain the same: the software selects the proposer after the method ballot; the proposal either requires three approvals from the other four citizens or takes effect directly; rejection returns the group to individual allocations.

\sectionbar{Unpaid expectation questions}
These questions appear before Block 1 Round 10 and before Block 2 Rounds 1 and \irbreplace{10}{5}. Participants enter their best estimates; answers do not affect payment.
\begin{enumerate}
  \item If group approval is required, how many points do you expect to earn? \hfill [0--60]
  \item If the selected person's decision takes effect directly, how many points do you expect to earn? \hfill [0--60]
  \item How many fund points do you expect the selected person to propose moving to their own payoff? \hfill [0--20]
\end{enumerate}

\sectionbar{Final questionnaire}
All paid decisions are complete before these items appear.

\begin{enumerate}
  \item How serious did the public-service situation in Block 1 seem?\\
  Not serious / Slightly serious / Moderately serious / Serious / Very serious
  \item Compared with Block 1, how did the public-service situation in Block 2 seem?\\
  Much worse / Somewhat worse / No change / Somewhat better / Much better
  \item After these rounds, how well can the group support public services when the selected person's proposal requires group approval?\\
  Very low / Low / Moderate / High / Very high
  \item After these rounds, how high is the risk that the usual limits on a leader's power will be seriously weakened?\\
  Very low / Low / Moderate / High / Very high
\end{enumerate}

For each statement below, participants use a slider from 1 (strongly disagree) to 7 (strongly agree):
\begin{enumerate}[resume]
  \item The current situation justifies allowing the selected person to act without group approval.
  \item More effective public services are worth giving up the approval vote in the current situation.
  \item Citizens should retain an approval vote over binding group decisions, even if agreement can fail.
  \item The leader should be allowed to require every citizen to put the same number of points in the public-services fund.
  \item I prefer requiring approval to letting a selected person implement an unwanted transfer.
  \item Removing the approval vote creates risks that outweigh the current gains.
  \item When public services work poorly, selected decision-makers need more freedom to act without further approval.
\end{enumerate}

\sectionbar{Optional background and predisposition questions}
The following twelve items appear only after the eleven process questions above and after all paid decisions. Each item is optional; participants may continue without answering it.
\begin{enumerate}
  \item \textbf{What is your age?}\\
  18--24 / 25--34 / 35--44 / 45--54 / 55--64 / 65 or older / Prefer not to answer
  \item \textbf{What is your gender?}\\
  Woman / Man / Another gender / Prefer not to answer
  \item \textbf{What is the highest level of education you have completed?}\\
  Basic or lower-secondary / General upper-secondary / Vocational qualification / Bachelor's or university of applied sciences degree / Master's / Doctoral / Other / Prefer not to answer
  \item \textbf{How easy or difficult is it for your household to make ends meet?}\\
  Very difficult / Difficult / Neither difficult nor easy / Easy / Very easy / Prefer not to answer
  \item \textbf{Which party did you vote for in the 2023 Finnish parliamentary election?}\\
  National Coalition Party / Finns Party / Social Democratic Party / Centre Party / Green League / Left Alliance / Swedish People's Party of Finland / Christian Democrats / Movement Now / Another party / I did not vote / I was not eligible to vote / Prefer not to answer
  \item \textbf{In politics, where would you place yourself?}\\
  0 (left) to 10 (right) / Prefer not to answer
  \item \textbf{Democracy may have problems, but it is better than any other form of government.}\\
  Strongly disagree / Somewhat disagree / Somewhat agree / Strongly agree / Prefer not to answer
  \item \textbf{On the whole, how satisfied are you with the way democracy works in Finland?}\\
  0 (completely dissatisfied) to 10 (completely satisfied) / Prefer not to answer
  \item \textbf{Would having a strong leader who does not have to bother with parliament and elections be a good or bad way of governing Finland?}\\
  Very bad / Fairly bad / Fairly good / Very good / Prefer not to answer
  \item \textbf{Citizens' opinions are taken into account in political decision making in Finland.}\\
  Strongly disagree / Disagree / Neither agree nor disagree / Agree / Strongly agree / Prefer not to answer
  \item \textbf{Generally speaking, would you say that most people can be trusted, or that you cannot be too careful?}\\
  0 (you cannot be too careful) to 10 (most people can be trusted) / Prefer not to answer
  \item \textbf{How willing are you to take risks in general?}\\
  0 (not at all willing) to 10 (very willing) / Prefer not to answer
\end{enumerate}

These are the complete planned questionnaire items. Minor corrections to language, translation, or screen formatting will be versioned. Any substantive new item, replacement, or construct will be submitted through the applicable ethics amendment or addendum process before use.

\sectionbar{Optional pretest feedback}
\irbadd{The first complete ten-person session offers this separate form after all paid decisions and main questionnaires, before the debrief and payment display. All participants see the same address. Five to ten responses are sought, and the researcher disables the invitation after this session. The form opens separately and does not contain the participant's oTree code or game responses. It stores aggregate rating counts and independent shuffled text lists, not linked individual questionnaires.}

\begin{infobox}[Invitation to give feedback]
\irbadd{Please help us check whether the experiment is clear and easy to use. These six questions are optional. Your payment and participation are unaffected if you skip them or criticize any part of the experiment. Your answers are anonymous and cannot be connected to your game answers. Please do not include names, study codes, seat numbers, or other identifying details. You may skip any question. When finished, close this tab and return to the experiment to continue to the final explanation and payment.}
\end{infobox}

\begin{enumerate}
  \item \irbadd{How clear were the instructions?}
  \item \irbadd{How clear was the difference between the two decision methods?}
  \item \irbadd{How clear was the explanation of your payment?}
\end{enumerate}
\irbadd{Questions 1--3 offer 1 (Very unclear), 2 (Unclear), 3 (Neither unclear nor clear), 4 (Clear), and 5 (Very clear). Each may be left unanswered.}
\begin{enumerate}[resume]
  \item \irbadd{Was anything confusing or difficult to do on the screens? Please describe.}
  \item \irbadd{What one change would most improve the experiment?}
  \item \irbadd{What do you think this study is trying to understand?}
\end{enumerate}
\irbadd{Questions 4--6 are optional text fields of at most 600 characters each.}

\sectionbar{Debriefing}
The study examines whether people use a group ballot to remove approval over a later binding decision and whether they subsequently vote to restore that approval. For some matching sets, public-service conditions improved and the two fund rates became equal; for others, the crisis and rate gap continued. Under both methods, the selected person could propose the same allocation and personal transfer; the difference was whether the other citizens could reject it.

No single choice is interpreted as evidence that a participant is corrupt, authoritarian, or opposed to democracy. Changing a choice is not treated as inconsistency. The study examines how institutional choices respond when public-service conditions and their material consequences change.

\end{document}
